\begin{problem}[5 points]{Complexity analysis}
  The following algorithm performs a linear search.
  State its worst-case running time and justify your answer.

  \begin{codebox}[python]{linear_search.py}
def locate(values, target):
    for index, value in enumerate(values):
        if value == target:
            return index
    return None
  \end{codebox}

  \begin{solution}
    In the worst case, the algorithm inspects every element exactly once.
    Its running time is therefore $\Theta(n)$ and its auxiliary space usage is $\Theta(1)$.
    The source and explanation are kept together following the general literate-programming principle described by \citet{knuth1984literate}.
  \end{solution}
\end{problem}
